\chapter{Introducción}
\label{ch:introduccion}

\section{Motivación}
% El reto de movilidad de la ZMG; dependencia del automóvil; el costo de una
% inversión mal asignada en transporte.
\placeholder

\section{El problema de circularidad en la planeación basada en datos}
% Los modelos que predicen dónde ya existe transporte ("tiene parada") premian el
% statu quo y no pueden identificar la necesidad no atendida.
\placeholder

\section{Preguntas de investigación}
\begin{description}
  \item[PI1 -- Demanda.] ¿Puede estimarse la demanda de transporte a resolución de
    \ageb solo a partir de datos de Nivel~1 (no de transporte)?
  \item[PI2 -- Diagnóstico.] ¿Un índice independiente de brecha de cobertura
    identifica zonas genuinamente desatendidas, y puede validarse fuera de muestra?
  \item[PI3 -- Generación.] ¿Pueden generarse y evaluarse corredores candidatos con
    base en necesidad genuina, no redundancia y eficiencia de demanda?
  \item[PI4 -- Transferibilidad.] ¿El marco se transfiere a otras zonas
    metropolitanas sin rediseño a la medida?
\end{description}

\section{Contribuciones}
\placeholder

\section{Estructura de la tesis}
% Un párrafo que mapea los capítulos a las ocho etapas.
\placeholder
